\section{R18: SPICE ephemeris representation and query}
\label{sec:r18-spice}

R18 constructs and queries actual CSPICE N0067 type 9 SPKs through
SpiceyPy 8.2.0. Every input state comes from the unchanged R10 native CPU
ORBIT-DYNAMICS-R2 forecast. No future observation is supplied to SPICE and no
model is refitted. This measures how well an SPK stores and reconstructs the
same forecast; it is not independent satellite-position validation or forward
physical propagation by SPICE. The underlying forecast still needs the separate
external-reference comparison.

\subsection{Same states, explicit time and frame conventions}
The initial declared grid covers G05 (MEO), C03 (GEO), C06 (IGSO) and Chandra
(HEO), from both January 2 and February 2, 2025, through seven days. These dates
were used in earlier releases. Before state evaluation, the protocol fixes
input intervals of 300, 900 and 1800 seconds crossed with polynomial degrees
7, 9 and 15: 72 cases. Every interval midpoint is compared with a native
state, every stored node is checked, and queries beyond either coverage end
must fail. The complete model, binary, interface, dependency and output hashes
are recorded in \path{review/r18_spice_protocol.json} and the comparison report.

For GPS epoch coordinate $g_0$ and model offset $t$, use
\begin{equation}
 T=(2444244.5-2451545.0)86400+g_0+51.184+t,
 \qquad E=\mathrm{UNITIM}(T,\mathrm{TDT},\mathrm{TDB}).
\end{equation}
Here $T$ and $E$ are TT and TDB seconds past J2000. The pinned NAIF LSK gives
$T=E-K\sin\epsilon$, $\epsilon=M+e_B\sin M$ and $M=M_0+M_1E$, hence
\begin{equation}
 \frac{dT}{dE}=1-K\cos\epsilon\,M_1(1+e_B\cos M).
\end{equation}
SPK velocities are multiplied by $dT/dE$ on writing and converted back on
reading. Metres become kilometres in the SPK. States are geometric relative
to Earth with aberration correction \code{NONE}. Model GCRS axes are transported
under an explicit identity-orientation convention to the SPICE J2000 label;
this does not validate a relativistic GCRS/ICRF coordinate transformation.

The tested TT round trips agree exactly at their floating-point inputs.
CSPICE's short LSK correction differs from geocentric ERFA by up to
15.320 microseconds in January and 22.406 microseconds in February; both SPK
nodes and queries use the same CSPICE convention. The stored model Julian TT
epoch differs by $-8.8215$ microseconds from the GPS-derived value. Absolute
ET doubles here have approximately 0.119 microsecond spacing. These numerical
interfaces remain finite precision despite lossless seed transport.

\subsection{Observed interpolation errors and a useful storage profile}
At the initial 300-second, degree-9 setting, the six GNSS cases have midpoint
maximum position differences of 0.455--1.082 mm. Chandra instead reaches
3.211 m in January and 1.667 m in February. Its coarse 900-second, degree-9
case reaches 8.716 km; the 1800-second, degree-15 February case reaches
45,580.644 km. All failures remain in the report. A higher polynomial degree
does not repair every undersampled trajectory or improve every boundary.

After those results, a separately declared adaptive follow-up tests Chandra
intervals of 30, 60 and 120 seconds with degrees 9 and 15. It is not a blind
confirmation. Its denser native grid preserves every earlier reference state
exactly. Degree-9 results are:
\begin{center}
\begin{tabular}{rrrr}
\toprule
Interval (s) & SPK bytes / seven days & January max (mm) & February max (mm)\\
\midrule
30 & 1,134,592 & 0.497 & 0.484\\
60 & 569,344 & 0.497 & 0.704\\
120 & 286,720 & 1.012 & 0.752\\
\bottomrule
\end{tabular}
\end{center}
Degree 15 sometimes worsens edge errors, reaching 14.626 mm in the February
30-second case. The measured starting profile is degree 9 with 300-second
GNSS nodes and 60-second Chandra nodes: at most 1.083 mm added position error
at tested midpoints, with 114 KiB per GNSS object and 556 KiB for Chandra.
The 120-second Chandra option uses 280 KiB and also stays below 1.013 mm at
the tested midpoints. These are sampled reconstruction results for these
trajectories, not continuous or future physical bounds.

The original experiment checks 72,648 nodes and 72,576 midpoints; the follow-up
adds 141,132 nodes and 141,120 midpoints. A fresh process reopens all 84 SPKs
and reproduces every reported maximum exactly. Every coverage rejection and
node-preservation check passes. Full protocols, arrays, SPKs and verification
reports are retained under \path{review/r18_spice*}.

\subsection{Cost, limits and physical meaning}
Five repeated warm batches of 2,016 shuffled queries at the original
300-second, degree-9 profile take median 0.457--0.463 s through the native
GNSS worker API, about 1.793 s for Chandra, and 0.00875--0.00935 s through the
SpiceyPy query API. SPICE is 49.4--51.9 times faster on these GNSS API cases
and 197.0--204.9 times faster on Chandra. It reuses a precomputed 116,736-byte
SPK, while the native path still executes RK4 steps from checkpoints and
crosses IPC/JSON. Precomputation is excluded and reported separately. Host load
is uncontrolled. These are storage/compute and interface measurements, not
equal arithmetic work, GPU throughput or a general product ranking.

For native forecast $x_N$, SPK reconstruction $x_S$ and external physical
reference $x_*$ at a common time and frame,
\begin{equation}
 \|x_S-x_*\|\leq\|x_S-x_N\|+\|x_N-x_*\|.
\end{equation}
This experiment measures the first term at specified samples. It enables
aTOMos forecasts to be distributed through established SPICE geometry
interfaces while preserving their measured accuracy to the chosen interpolation
budget. It does not improve the second term or establish a physical accuracy
advantage over NASA software. Detailed equations, chronology and reproduction
commands are in \path{formal/SPICE_EPHEMERIS_COMPARISON.md}.
The primary interface definitions are NASA NAIF's
\href{https://naif.jpl.nasa.gov/pub/naif/toolkit_docs/C/cspice/spkw09_c.html}{SPK type 9}
and \href{https://naif.jpl.nasa.gov/pub/naif/toolkit_docs/C/cspice/unitim_c.html}{uniform-time conversion}
documentation; the numerical results above are local measurements.
